\section{Model and numerical methods}
\label{sec:methods}

\subsection{Hamiltonian and isolated-cell projection}

The Hamiltonian of the pentamer chain is
\begin{align}
H={}&J_1\sum_n\left(\bm S_{n,2}\!\cdot\!\bm S_{n,3}
 +\bm S_{n,3}\!\cdot\!\bm S_{n,4}\right) \notag\\
&+J_2\sum_n\left(\bm S_{n,1}\!\cdot\!\bm S_{n,2}
 +\bm S_{n,4}\!\cdot\!\bm S_{n,5}\right) \notag\\
&+J_3\sum_n\bm S_{n,5}\!\cdot\!\bm S_{n+1,1},
\label{eq:hamiltonian}
\end{align}
where $n$ labels the five-site unit cells and $\bm S_{n,a}$ is the spin-$1/2$
operator at site $a$ of cell $n$.
The intra-cell bonds are the two outermost bonds at $J_2$ and the two central
bonds at $J_1$; a single bond $J_3$ connects neighboring cells.
We work in units of $J_1$ and set $\alpha=J_2/J_1$ and $\beta=J_3/J_1$ throughout.

The isolated pentamer has a 32-dimensional Hilbert space.
We diagonalize it completely and resolve each degenerate energy subspace by total
spin $S$ and by reflection parity about the cell center.
Spectral weights are evaluated with projectors onto complete multiplets rather
than with individual eigenvectors, whose orientation inside a degenerate subspace
is arbitrary.
This matters here because the multiplet structure, not merely the energy
ordering, controls the dynamical selection rules derived below.

If $P$ projects onto the ground-state doublet, the local spin operators act
within that subspace as
\begin{equation}
P\bm S_{n,a}P=c_a\bm T_n,
\label{eq:projection}
\end{equation}
where $\bm T_n$ is an effective spin-$1/2$ operator and the $c_a$ are
cell-dependent projection coefficients.
To leading order in the inter-cell coupling the chain is then described by
\begin{equation}
\Heff=NE_5+\Jeff\sum_n\bm T_n\!\cdot\!\bm T_{n+1}+O(\beta^2),
\qquad \Jeff=\beta c_1c_5,
\label{eq:effective}
\end{equation}
with $N$ the number of cells and $E_5$ the isolated-cell ground-state energy.
Only the end sites $a=1$ and $5$ enter $\Jeff$, because they are the sites touched
by the inter-cell bond.
Reflection about the cell center maps site 1 onto site 5 and leaves an even
doublet invariant, so $c_5=c_1$ and $\Jeff=\beta c_1^2>0$: an even ground doublet
always produces an antiferromagnetic effective chain.
The parity of the cell ground state thus fixes the sign of the effective coupling,
and its size is set by how much of the doublet's spin resides on the cell ends.
The same coefficients determine what a scattering probe sees.
Because each effective spin is spread over the cell with amplitudes $c_a$, a
probe at physical momentum $q$ couples to it with the cell shape factor
\begin{equation}
F(q)=\Bigl|\sum_{a=1}^{5}c_a e^{iqa}\Bigr|^2,
\label{eq:shape}
\end{equation}
which folds the effective-chain response back into the physical Brillouin zone
and fixes the momenta at which it appears.
We test the projection in two ways.
Exact diagonalization of two coupled pentamers measures the departure of the
exact ground-state energy from the projected form, and finite-chain calculations
with $N=8,12,16$ cells test whether the same leading power persists in the
extended system.
Figure~\ref{fig:model} summarizes the geometry, the isolated-cell multiplet
structure, and the projection that defines the effective chain.

\begin{figure*}[t]
  \includegraphics[width=\textwidth]{fig1_model_projection.png}
  \caption{Model and low-energy projection.
  The five-site unit cell defines the couplings $J_1$, $J_2$, and $J_3$.
  Isolated-pentamer multiplets are resolved by total spin and reflection parity,
  and the even-parity ground-state doublet is projected onto an effective
  spin-$1/2$ degree of freedom that mediates the inter-cell exchange.}
  \label{fig:model}
\end{figure*}

\subsection{Finite-chain calculations and convergence criteria}
\label{sec:convergence}

Finite chains use open boundaries and $L=5N$ physical spins.
Ground states are obtained in fixed total-$S^z$ sectors with density-matrix
renormalization group calculations in a matrix-product-state
representation~\cite{white1992,schollwoeck2005,schollwoeck2011}.
Sector-specific initial states are used rather than assuming that a single
initialization samples every low-energy solution.

The positive-region data set contains 105 converged calculations at
$L=40,60,80$.
For negative $\alpha$, we initialize the $S^z=0$ sector from one of two distinct
N\'eel-type product states, and the $S^z=1$ sector from a single fixed product
state.
A calculation is retained only if the variance per site satisfies
$\mathrm{Var}(H)/L\le5\times10^{-8}$, the final truncation error satisfies
$\epsilon_{\mathrm{trunc}}\le10^{-9}$, the measured magnetization satisfies
$|S^z-S^z_{\mathrm{target}}|\le10^{-8}$, and the sector gap
$E_0(S^z=1)-E_0(S^z=0)$ is no smaller than $-10^{-8}$, so that the $S^z=0$ sector
holds the ground state; each calculation uses at most 64 sweeps.
The two $S^z=0$ initializations probe sensitivity to the starting state.

Every calculation was re-evaluated before use, and only calculations that pass
these criteria enter the analysis below.
These criteria test the variational state at a given solver setting; the
sensitivity to the eigensolver itself is examined separately in
Sec.~\ref{sec:krylov}.

\subsection{Entanglement, correlation, and momentum diagnostics}
\label{sec:diagnostics}

For an open critical chain the bond entropy is compared with the conformal form
\begin{equation}
S(b)=a_\nu+\frac{\ceff}{6}
\ln\left[\frac{2L}{\pi}\sin\left(\frac{\pi b}{L}\right)\right],
\label{eq:entropy}
\end{equation}
where the intercept $a_\nu$ depends on the bond class $\nu=b\bmod 5$ or $b\bmod 10$
and a common slope defines the descriptive coefficient $\ceff$.
Fits use central windows $fL\le b\le(1-f)L$ with $f=0.10$, 0.15, 0.20, and 0.25.
Equation~\eqref{eq:entropy} is the standard conformal result for the entanglement
entropy of an open critical
chain~\cite{calabrese2004,giamarchi2003}; because the five-site unit cell makes
bond positions inequivalent, we treat the intercept as bond-class dependent and
refer to the common slope as the effective entropy coefficient.

Correlation envelopes, finite-size sector gaps, the static structure factor, and
connected bond--bond correlations provide complementary diagnostics.
Momentum-resolved information is obtained from the static structure factor
$S(q)=\sum_{r}e^{iqr}\avg{\bm S_0\!\cdot\!\bm S_r}$ and from the dynamical
structure factor defined in Sec.~\ref{sec:dsf}.

The negative-$\alpha$ scan tracks the momentum $\qpeak$ at which $S(q)$ is
largest and the sharpness of that peak,
\begin{equation}
R_L=1-\frac{S(\qpeak-\delta)+S(\qpeak+\delta)}{2S(\qpeak)},
\label{eq:ratio}
\end{equation}
where $\delta$ is the momentum resolution $2\pi$ divided by the number of sites
in the central window.
A peak that sharpens with length, $\Delta R=R_{80}-R_{60}>0$, signals a developing
singularity, whereas a peak that broadens points away from one.
A candidate boundary is retained only when an adjacent coupling interval shows
both a momentum reconstruction and a sign change of $\Delta R$.
A momentum change without the corresponding $\Delta R$ condition is reported as
a finite-size reconstruction, not as a phase boundary.

\subsection{Dynamical structure factor}
\label{sec:dsf}

We compute the dynamical structure factor
\begin{equation}
S^{zz}(q,\omega)=\sum_{m}\left|\avg{m|S^z_q|0}\right|^2
\delta\!\left(\omega-E_m+E_0\right)
\label{eq:dsf}
\end{equation}
by two independent matrix-product-state routes.
Correction-vector DMRG gives a fixed-frequency spectrum at $q=\pi$ by solving the
linear system associated with the resolvent of $H$ at complex frequency
$\omega+\mathrm{i}\eta$~\cite{kuehner1999,jeckelmann2002,nocera2016,nocera2018rixs}.
Real-time evolution followed by Fourier transformation provides the full momentum
dependence and an independent line-shape check~\cite{vidal2004,paeckel2019}.
As a validation of the Hamiltonian, the momentum source, the broadening, and the
spectral normalization, the correction-vector implementation reproduces the
complete $L=10$ exact spectrum with a relative $L_2$ error of $9\times10^{-12}$.

The production full-momentum calculations use $L=40$, $t_{\max}=30$, 121 momentum
points, 301 frequency points, and paired time steps $\Delta t=0.025$ and
$0.0125$.
The raw Fourier matrices, including negative ringing from the finite time window,
are retained and tested rather than clipped, so that the convergence assessment
rests on the full spectral weight.

\subsection{Krylov-subspace stability control}
\label{sec:krylov}

At $(\alpha,\beta)=(-0.10,0.25)$ an earlier $L=80$ analysis gave an entropy
coefficient well below the values near unity found elsewhere
(Sec.~\ref{sec:saved-state}).
Because this point was also the most sensitive to the eigensolver configuration,
we performed a final $L=80$ control in which all physical and DMRG settings were held fixed while the
Krylov dimension was increased from $k=8$ to $k=12$.
Both initializations used 64 sweeps, a configured maximum bond dimension of 1200,
a cutoff of $10^{-11}$, and the same 401-point momentum grid.
Here $k$ is the Krylov-subspace dimension, distinct from the MPS bond dimension
$\chi$.

The $k=12$ calculation satisfied the same convergence criteria listed in
Sec.~\ref{sec:convergence}.
Its variance per site was $1.9\times10^{-11}$, and the two sector gaps were
$7.2772\times10^{-5}$, agreeing to five significant figures.
Across 172 matched observables, the $k=8\rightarrow12$ differences were bounded by
$7.5\times10^{-12}$ for the spin gap, $8.4\times10^{-8}$ for the dominant
ratio, and $5.7\times10^{-8}$ for the entropy, while $\qpeak/\pi=1$ was
unchanged.
The largest retained link dimension was $264$--$265$ in the $S^z=0$ sector and
$186$--$190$ in $S^z=1$, nearly insensitive to the solver change and well below
the configured maximum of 1200, so truncation at this point is set by the
$10^{-11}$ cutoff rather than by the bond-dimension cap.

The $L=40$ and $L=60$ members of the negative-$\alpha$ size sequence in
Sec.~\ref{sec:saved-state} were obtained with the earlier setting $k=3$; the
quantitative statements there rest on the controlled $L=80$ endpoint.
